\thispagestyle{plain}

% ── HEADER ────────────────────────────────────────────────────────────────────
\begin{center}
  {\large\bfseries General Directorate of Technological Studies}\\
  {\large\bfseries Higher Institute of Technological Studies of Beja}\\
  {\large\bfseries Department of Information Technology}
\end{center}

\vspace{1.2cm}

% ─ BOXED ABSTRACT TABLE ──────────────────────────────────────────────────────
\begin{center}
\renewcommand{\arraystretch}{1.3}
\begin{tabular}{|p{14.5cm}|}
\hline
\multicolumn{1}{|c|}{\textbf{Graduation Project in Development of Embedded and Mobile Systems}} \\
\hline
\textbf{Made by:} Fedi Amdouni \hfill \textbf{Batch:} 2025-2026 \newline
 \hspace*{1.86cm} Firas Ouesleti \\
\hline
\textbf{Title:} Dev and Deployment of an automated real-time monitoring system \\
\hline
\textbf{Defence:} 18/06/2026 \\
\hline
\textbf{Summary:}
\vspace{0.1cm}

In modern manufacturing, accurate production monitoring is essential for optimizing efficiency. However, many facilities face unreliable traceability of machine stoppages and their root causes, hindering precise Overall Equipment Effectiveness (OEE) calculation and performance analysis. This project addresses this challenge by designing and deploying an automated, real-time industrial monitoring system. Built on a multi-layer architecture, the solution integrates PLCs via Modbus TCP for data acquisition, Node-RED for real-time processing, NestJS for backend management, and a custom HMI for operator interaction. The system automatically detects stoppages, enables real-time cause logging, and centralizes production data in a MySQL database, while analytical dashboards provide dynamic KPI visualization. Implementation results demonstrate complete machine event traceability, automatic OEE calculation per shift and day, and a significant reduction in manual entry errors and undetected production losses. The system successfully enhances production transparency and decision-making, with future work planned for scaling across additional production lines.
\vspace{0.2cm}
\\
\hline
\textbf{Keywords:} Industrial monitoring, OEE, Modbus TCP, PLC, Node-RED, NestJS, HMI, machine stoppages, MySQL database \\
\hline\hline  \hline  
\textbf{Résumé:}
\vspace{0.1cm}

Dans l'industrie moderne, une surveillance précise de la production est essentielle pour optimiser l'efficacité. Cependant, de nombreuses installations font face à une traçabilité peu fiable des arrêts machine et de leurs causes profondes, entravant le calcul précis de l'Efficacité Globale des Équipements (OEE) et l'analyse des performances. Ce projet relève ce défi en concevant et déployant un système de surveillance industrielle automatisé et en temps réel. Construite sur une architecture multi-couches, la solution intègre des automates programmables (PLC) via Modbus TCP pour l'acquisition des données, Node-RED pour le traitement en temps réel, NestJS pour la gestion backend, et une interface homme-machine (IHM) personnalisée pour l'interaction des opérateurs. Le système détecte automatiquement les arrêts, permet l'enregistrement en temps réel des causes, et centralise les données de production dans une base de données MySQL, tandis que des tableaux de bord analytiques fournissent une visualisation dynamique des indicateurs de performance (KPI). Les résultats de l'implémentation démontrent une traçabilité complète des événements machine, un calcul automatique de l'OEE par équipe et par jour, et une réduction significative des erreurs de saisie manuelle et des pertes de production non détectées. Le système améliore avec succès la transparence de la production et la prise de décision, avec des travaux futurs prévus pour l'extension à d'autres lignes de production.
\vspace{0.2cm}
\\
\hline
\textbf{Mots-clés:} Surveillance industrielle, OEE, Modbus TCP, PLC, Node-RED, NestJS, IHM, arrêts machine, base de données MySQL \\
\hline
\end{tabular}
\end{center}

\vfill